\section{Related Work}
\label{sec:related}

\para{Chain-of-thought prompting.}
\citet{wei2022chain} showed that including step-by-step reasoning exemplars in prompts dramatically improves multi-step reasoning in LLMs, achieving state-of-the-art results on benchmarks like \gsm with PaLM 540B. \citet{kojima2022large} extended this to the zero-shot setting, demonstrating that simply appending ``Let's think step by step'' triggers effective reasoning without exemplars. These approaches place all reasoning \emph{before} the final answer. \citet{yao2023tree} generalized CoT with tree-structured exploration of reasoning paths. Our work differs in that we interleave reasoning \emph{within} the response at sentence granularity rather than concentrating it before the answer.

\para{Pause and filler tokens.}
\citet{goyal2023think} proposed appending learnable \texttt{<pause>} tokens to give models extra computation before generating output, improving performance on 8 of 9 tasks with a 1B-parameter model. Critically, they found that pause-pretraining is essential---retrofitting pauses onto standard models yields limited benefit, and larger models benefit more. \citet{pfau2024lets} showed that even meaningless filler tokens (dots) enable transformers to solve tasks they otherwise cannot, recovering approximately 90\% of CoT benefit on parallelizable tasks. Both approaches require training modifications. We test whether prompting alone can achieve similar benefits, avoiding the need for model retraining.

\para{Latent and interleaved reasoning.}
Quiet-STaR~\cite{zelikman2024quietstar} trains models to generate internal rationales at every token position using REINFORCE, with a mixing head to interpolate between base and post-thought predictions. The Self-Notes framework~\cite{lanchantin2023selfnotes} showed that reasoning interleaved near relevant content outperforms post-context reasoning (scratchpad), with interleaved notes achieving 85.0\% accuracy on out-of-distribution algorithmic tasks versus 11.6\% for scratchpad. Both require fine-tuning. Our approach operates at the same conceptual level as Self-Notes---interleaving reasoning near the content it verifies---but uses prompting rather than training, and operates at sentence granularity rather than arbitrary positions.

\para{Implicit and compressed reasoning.}
Several works explore internalizing explicit reasoning into model representations. \citet{deng2023implicit} distill CoT into hidden-state computation via knowledge distillation, while \citet{deng2024explicit} gradually remove CoT tokens during fine-tuning to internalize reasoning. \citet{hao2024coconut} propose reasoning in continuous latent space by feeding hidden states back as inputs. \citet{cheng2024compressed} compress full reasoning chains into dense ``contemplation tokens.'' These approaches reduce or eliminate the token overhead of explicit reasoning. Our work highlights why this matters: we show that the token budget consumed by explicit thinking tokens is the primary limitation of prompt-based interleaved reasoning.

\para{Iterative refinement.}
Self-Refine~\cite{madaan2023selfrefine} uses iterative generate-critique-refine cycles to improve output quality, achieving 20\% improvement on average. \citet{nye2021scratchpads} demonstrated early on that training models to emit intermediate computation steps improves multi-step tasks. STaR~\cite{zelikman2022star} bootstraps reasoning by having models learn from their own generated rationales. Unlike iterative refinement, which operates as a post-hoc correction process, \sentthink integrates deliberation into the initial generation process, verifying claims as they are produced rather than after the fact.
